% Appendix proof module. Compatible with scalar_main.tex but mostly standalone.
\section{Scalar theory: assumptions and complete derivations}\label{sc:proofs}

\subsection{Deterministic and iid quadratic products}
For \(\ell(x)=hx^2/2\), deterministic gradient descent satisfies
\(x_t=(1-\eta h)^t x_0\). If \(h>0\), convergence to zero from every
finite start occurs precisely when \(0<\eta h<2\). At \(\eta h=2\),
nonzero starts follow a two-cycle rather than converge, and at
\(\eta h>2\) their magnitudes diverge. The mean, mean-square, logarithmic,
and harmonic multiplier criteria coincide in this deterministic case.
This calculation does not describe stationary fluctuations because
there is no sampled forcing and no mechanism returning an expanding
linear orbit to a bounded scale.

For iid batches, common-noise differences satisfy
\(\Delta_t/\Delta_0=\prod_{j=1}^t A_j\), regardless of additive forcing.
The logarithm is a sum of iid integrable variables, so the strong law
gives the almost-sure exponent in Proposition~\ref{sc:ordering}.
Independence also gives the exact \(p\)-moment identity, for every real
\(p\) for which the moment is defined. For negative \(p\), this identity
describes inverse-distance growth; its contraction means repulsion from
zero and must not be called contraction toward zero.

Jensen applied to \(\exp(-\log M)\) gives
\(\log\mathbb E M^{-1}\ge-\mathbb E\log M\). Jensen applied to
\(\exp(\log M)\) gives \(\mathbb E\log M\le\log\mathbb E M\).
Cauchy--Schwarz gives \(\mathbb E M\le(\mathbb E M^2)^{1/2}\).
These are precisely \eqref{sc:ladder}. If \(M\) is nonconstant, the
first Jensen inequality is strict whenever both sides are finite, so
\(H=1\) then also implies \(\gamma>0\).

For example, let \(M=1/2\) with probability \(1/5\) and \(M=2\) with
probability \(4/5\). Then \(\gamma=(3/5)\log2>0\),
\(m(-1)=4/5\), \(H=5/4>1\), and \(m(-2)=1\).
Changing the contracting probability to \(2/5\) gives positive
\(\gamma=(1/5)\log2\) while \(m(-1)=11/10>1\), hence \(H<1\).
The separated thresholds therefore occur already for two-point curvature
laws, not only for estimated high-dimensional diagnostics.

\subsection{The affine baseline and why outer return matters}
The following restricted setting suffices for an exact comparison and
avoids unspoken nondegeneracy exceptions. Let \((A_t,N_t)\) be iid,
\(A_t\) take finitely many nonzero real values, and, conditionally on
\(A_t\), let \(N_t\) be centered Gaussian with variance in
\([v_-,v_+]\subset(0,\infty)\). Multipliers and noise may be dependent.
Consider
\begin{equation}
 X_{t+1}=A_{t+1}X_t+r(X_t)+N_{t+1},\qquad \sup_x|r(x)|\le R<\infty,
 \label{sc:bounded-remainder}
\end{equation}
where \(r\) is continuous.

\begin{proposition}[Outer logarithmic drift separates recurrence and escape]
\label{sc:outer-sign}
If \(\mathbb E\log|A|<0\), \eqref{sc:bounded-remainder} has a unique
positive Harris invariant probability and converges geometrically in a
weighted total-variation norm with weight \(1+|x|^p\) for some
\(0<p\le1\). If \(\mathbb E\log|A|>0\), it has no invariant probability.
No conclusion at equality is asserted.
\end{proposition}
\begin{proof}
In the negative case, finiteness and nonvanishing of the multiplier
support imply that \(q(p)=\mathbb E|A|^p\) is differentiable at zero,
with \(q(0)=1\) and \(q'(0)<0\). Choose \(0<p\le1\) with \(q(p)<1\).
Subadditivity gives
\[
 P(1+|x|^p)\le q(p)(1+|x|^p)+1-q(p)+R^p+\mathbb E|N|^p.
\]
The transition density is a finite mixture of strictly positive Gaussian
densities, continuous jointly in the starting and target state. On any
two fixed compact intervals it has a positive lower bound. Every compact
set is therefore small, the chain is Lebesgue irreducible and aperiodic,
and the geometric Foster--Harris theorem applies
\citep{hairermattingly2011}. To localize the drift, pick
\(q(p)<q'<1\) and absorb its additive constant inside a sufficiently
large level set of \(1+|x|^p\).

For the positive case write \(a_j=|A_j|\),
\(P_n=\prod_{j=1}^n a_j\), and \(d=\mathbb E\log a_1>0\).
The reverse triangle inequality, iterated, gives
\[
 \frac{|X_n|}{P_n}\ge |X_0|-
 \sum_{j=1}^n\frac{R+|N_j|}{P_j}.
\]
Fix \(0<c<d\). The strong law implies
\(\log P_n-cn\to\infty\), so
\(\inf_{n\ge0}P_ne^{-cn}>0\) almost surely. There is a deterministic
\(K>0\) such that
\(E_K=\{P_n\ge Ke^{cn}\text{ for all }n\ge0\}\) has positive
probability. Conditionally on the entire multiplier sequence,
\(\mathbb E(|N_j|\mid(A_k)_{k\ge1})\le C=\sqrt{2v_+/\pi}\).
On \(E_K\), therefore,
\[
 \mathbb E\!\left[\sum_{j\ge1}\frac{R+|N_j|}{P_j}
       \,\middle|\,(A_k)_{k\ge1}\right]
 \le D:=\frac{R+C}{K(e^c-1)}.
\]
For \(|X_0|\ge4D\), conditional Markov inequality gives probability at
least \(1/2\) that the weighted noise sum is at most \(|X_0|/2\).
Consequently the probability of \(|X_n|\to\infty\) from every such
large start is at least \(\mathbb P(E_K)/2>0\). The everywhere positive
one-step density reaches this region from every start with positive
probability. An invariant law \(\pi\) would thus give positive
stationary probability \(\delta\) to escape. For any \(M<\infty\),
Fatou's lemma and stationarity imply
\[
 \delta\le\mathbb E_\pi\liminf_n\mathbf1_{\{|X_n|>M\}}
 \le\liminf_n\mathbb P_\pi(|X_n|>M)=\pi(|x|>M).
\]
Letting \(M\to\infty\) contradicts that \(\pi\) is a probability.
\end{proof}

For the affine case \(r=0\), negative logarithmic drift also gives the
classical perpetuity representation
\(X\overset d=\sum_{n\ge1}(A_1\cdots A_{n-1})N_n\).
The products decay exponentially almost surely and Gaussian innovations
grow subexponentially, so the series converges absolutely. Splitting off
the first pair proves the distributional recursion. Under the additional
Kesten--Goldie assumptions (a root \(\mathbb E|A|^\kappa=1\),
\(0<\mathbb E|A|^\kappa\log|A|<\infty\), nonarithmetic
\(\log|A|\), the required \(\kappa\)-moment of the innovation, and no
deterministic common fixed point), the stationary tails have power
index \(\kappa\). This last tail theorem is classical
\citep{kesten1973random,goldie1991implicit}; the theorem above proves the
existence/nonexistence assertion needed here directly. Gaussian forcing
excludes absorbing point masses and common deterministic fixed points.

\subsection{Proof of the reflected-log central law}
Let \(Z_j=-\log M_j\), \(S_n=\sum_{j=1}^nZ_j\), \(S_0=0\), and
\(W=\log(r_*/R)\). Then
\[
 W_{t+1}=\max\{0,W_t+Z_{t+1}\},\qquad \mathbb EZ_1=-\gamma<0.
\]
The strong law implies \(S_n\to-\infty\), hence
\(W_*:=\sup_{n\ge0}S_n<\infty\) almost surely. Splitting the supremum
at its first increment yields
\(W_*\overset d=\max\{0,Z_1+W_*'\}\), with an independent copy
\(W_*'\). This proves invariance. To see uniqueness and attraction,
couple two finite starting values with the same increments. The larger
process reaches zero almost surely: before its first visit its value is
its starting value plus \(S_n\), which cannot stay positive forever.
The smaller process is zero at that time by monotonicity, and the two
processes then agree. Couple an arbitrary finite start to an independent
stationary start this way. The coupling time is finite almost surely,
so its tail tends to zero and total variation convergence follows.

Put \(\phi(q)=\mathbb E e^{qZ}=m(-q)\). At the positive root
\(\phi(\alpha)=1\), define the tilted increment law
\(F_\alpha(dz)=e^{\alpha z}F(dz)\). Its positive mean is
\(\mu_\alpha=\mathbb E[Ze^{\alpha Z}]\). For
\(T_u=\inf\{n\ge1:S_n>u\}\), the likelihood identity, first at each
event \(\{T_u=n\}\) and then summed over \(n\), gives
\begin{align*}
 \mathbb P(W_*>u)
 &=\sum_{n\ge1}\mathbb E_\alpha
       [e^{-\alpha S_n}\mathbf1_{\{T_u=n\}}]\\
 &=e^{-\alpha u}\mathbb E_\alpha e^{-\alpha O_u},
 \qquad O_u=S_{T_u}-u.
\end{align*}
Under the tilted law \(T_u<\infty\) almost surely, since its drift is
positive. The nonarithmetic overshoot renewal theorem, with the stated
exponential-neighborhood and finite-mean assumptions, gives
\(O_u\Rightarrow O_\infty\), where \(O_\infty\) is finite and
nonnegative. The bounded continuous function \(e^{-\alpha o}\) yields
\(\mathbb E_\alpha e^{-\alpha O_u}\to
C:=\mathbb E_\alpha e^{-\alpha O_\infty}\in(0,1]\).
This is the classical Cram\'er--Lundberg asymptotic; its random-walk
renewal formulation is also described by
\citet{blanchet2006corrected}.
Substitution of \(u=\log(r_*/r)\) proves the small-ball law.
Strict versus weak inequalities have the same asymptotic by squeezing
between radii \(r\) and \((1+\epsilon)r\) and then letting
\(\epsilon\downarrow0\).

Here is a proof of the concrete density condition in the theorem.
Suppose \(Z=-\log M\) has a bounded continuous density \(f\), supported
in \([-b_-,b_+]\), with \(b_+,b_->0\). Let \(\tau_+\) be its first
strictly positive ladder epoch and define the defective ladder-height
measure \(G(D)=\mathbb P(\tau_+<\infty,S_{\tau_+}\in D)\) on
\((0,b_+]\). Write \(p=G((0,b_+])\). Since \(S_n\to-\infty\), the
walk has finitely many strict upper records almost surely. The strong
Markov property would produce infinitely many records if \(p=1\), so
\(p<1\). Decomposing the supremum by its number of records gives
\[
 \mathcal L(W_*)=(1-p)\sum_{n\ge0}G^{*n}.
\]
The ladder measure has density
\[
 g_G(x)=\mathbf1_{\{0<x\le b_+\}}
   \sum_{n\ge0}\mathbb E[
      \mathbf1_{\{\tau_+>n\}}f(x-S_n)].
\]
The summand is zero unless \(-b_+\le S_n\le0\). A bounded-increment
random walk with negative mean has finite expected occupation of this
compact interval: for large \(n\), the event \(S_n\ge-b_+\) is an
exponentially unlikely deviation from its negative mean, by Hoeffding's
inequality. Thus \(g_G\) is bounded. Dominated convergence using the
same finite occupation measure makes it continuous on \((0,b_+)\),
with finite one-sided endpoint limits. It is consequently directly
Riemann integrable and compactly supported.

Exponential tilting at the stopped ladder epoch gives
\(\widehat G(dx)=e^{\alpha x}G(dx)\), which is a probability: under
the tilted positive-drift walk the first positive ladder epoch is finite
almost surely. Its mean \(\widehat\mu=\int x\widehat G(dx)\) lies
in \((0,b_+]\), and it has a bounded, directly Riemann integrable
density \(\widehat g\). Let
\(\widehat U=\sum_{n\ge0}\widehat G^{*n}\). The positive-increment
key renewal theorem applied to \(\widehat g\) gives
\[
 \widehat u(w):=\int\widehat g(w-y)\widehat U(dy)
       \longrightarrow\frac1{\widehat\mu}.
\]
This convolution is a density of \(\widehat U-\delta_0\), since
\(\widehat U=\delta_0+\widehat G*\widehat U\). Undoing the tilt
in the supremum decomposition therefore gives a density on \((0,\infty)\),
\[
 g_{W_*}(w)=(1-p)e^{-\alpha w}\widehat u(w)
 \sim\frac{1-p}{\widehat\mu}e^{-\alpha w}.
\]
Its integral tail identifies the earlier constant as
\(C=(1-p)/(\alpha\widehat\mu)\). Finally
\(p_R(r)=g_{W_*}(\log(r_*/r))/r\), proving the density assertion.
The atom \(1-p\) at \(W_*=0\), equivalently at \(R=r_*\), is retained.
This argument derives the local density from the positive-increment
renewal theorem and does not differentiate a CDF asymptotic.

For the harmonic comparison, let \(\psi(q)=\log\phi(q)\).
It is strictly convex for a nonconstant gain, with \(\psi(0)=0\) and
\(\psi'(0)=-\gamma<0\). Its unique positive zero is \(\alpha\).
It is negative on \((0,\alpha)\) and positive beyond \(\alpha\)
within its finite domain. Since \(\psi(1)=-\log H\), this proves
\eqref{sc:shape-ladder}. Existence of the root is a substantive assumption:
if \(M>1\) almost surely then no positive inverse root exists.

\subsection{A solvable return model with an exact density}
\begin{proposition}[Exponential log excursions]\label{sc:solvable}
Let \(U\sim\operatorname{Exp}(b)\), \(V\sim\operatorname{Exp}(c)\) be
independent, \(b>c>0\), and \(M=e^{V-U}\). In
\eqref{sc:saturation}, set \(\alpha=b-c\) and \(\rho=c/b\). Then
\begin{equation}
 \pi_R=(1-\rho)\delta_{r_*}
 +\rho\,\frac\alpha{r_*}\left(\frac r{r_*}\right)^{\alpha-1}
             \mathbf1_{(0,r_*)}(r)\,dr.
 \label{sc:exact-density}
\end{equation}
The local logarithmic drift is \(1/c-1/b>0\) and
\(m(-q)=bc/[(b-q)(c+q)]\) for \(0\le q<b\).
Thus \(m(-\alpha)=1\), and, when \(b>1\), \(H>1\) precisely when
\(\alpha>1\).
\end{proposition}
\begin{proof}
For the reflected coordinate, propose
\(W\sim(1-\rho)\delta_0+\rho\operatorname{Exp}(\alpha)\), independent
of fresh \(U,V\). Its Laplace transform is
\[
 \mathbb E e^{-sW}=1-\rho+\rho\frac\alpha{\alpha+s}
   =\frac{\alpha(b+s)}{b(\alpha+s)}.
\]
Multiplying by the Laplace transform \(b/(b+s)\) of \(U\) shows
\(W+U\sim\operatorname{Exp}(\alpha)\). Therefore, for \(w\ge0\),
\[
 \mathbb P((W+U-V)_+>w)
 =e^{-\alpha w}\mathbb E e^{-\alpha V}
 =\rho e^{-\alpha w}.
\]
This is the proposed invariant law; uniqueness follows from the preceding
coupling argument. Transforming back gives \eqref{sc:exact-density}.
Independence gives the displayed moment formula, whose nonzero root is
\(b-c\). Finally
\(H=(b-1)(c+1)/(bc)\), so
\(H-1=(\alpha-1)/(bc)\).
\end{proof}
The boundary atom is real and must be retained in plots or simulations.
A fair independent sign gives a symmetric invariant coordinate with
half the mass on each side, but its two boundary atoms are not a theorem
about two smooth interior modes in neural-network coordinates.

\subsection{Finite-mean reinjection: a useful generalization}
The hard cap can be replaced by independent excursions. The following
version uses bounded increments and entries to make the renewal averaging
step transparent. Suppose \(Y=\log M\) has a bounded continuous density
with bounded support, \(\mathbb EY>0\), \(\mathbb P(Y<0)>0\), and
\(\mathbb E e^{-\alpha Y}=1\) for \(\alpha>0\). Let the independent
entry \(Z\) lie in \([-L,-\ell]\), with \(0<\ell\le L<\infty\).
Run \(Q_n=Z+\sum_{j=1}^nY_j\) until
\(\tau=\inf\{n\ge0:Q_n\ge0\}\); discard the exit and restart at a
fresh independent \(Z\). Then
\begin{equation}
 \pi_Q(D)=\frac{\mathbb E\sum_{n=0}^{\tau-1}
                  \mathbf1_{\{Q_n\in D\}}}{\mathbb E\tau}
 \label{sc:occupation}
\end{equation}
is invariant. To see \(\mathbb E\tau<\infty\), write
\(\mu=\mathbb EY>0\). Since \(\{\tau>n\}\) implies
\(\sum_{j=1}^nY_j<L\), bounded-increment Hoeffding bounds make the
tail summable for \(n>2L/\mu\).
For bounded test functions, the expected one-step change summed over a
cycle telescopes to the next fresh entry minus the present entry; its
expectation is zero. This proves \eqref{sc:occupation} is invariant.

Let \(m_\alpha=-\mathbb E[Y e^{-\alpha Y}]>0\) and
\(D_\alpha=\mathbb E[e^{-\alpha Z}-e^{-\alpha Q_\tau}]>0\).
For every fixed \(h>0\), the renewal calculation gives
\begin{equation}
 \pi_Q((x,x+h])\sim
 \frac{D_\alpha}{\alpha m_\alpha\mathbb E\tau}
 (e^{\alpha h}-1)e^{\alpha x},\qquad x\to-\infty.
 \label{sc:regeneration-asymptotic}
\end{equation}
Here is the calculation. For the increment law \(F\), write
\(U=\sum_{n\ge0}F^{*n}\). Tilting by \(e^{-\alpha y}\) yields a
probability \(F_\alpha\) with negative drift \(-m_\alpha\) and
\(U(dt)=e^{\alpha t}U_\alpha(dt)\). The nonarithmetic random-walk
renewal theorem, applied on fixed intervals and then to continuous test
functions on those intervals, gives
\[
 U_z((x,x+h])\sim
 e^{-\alpha z}\frac{e^{\alpha h}-1}{\alpha m_\alpha}e^{\alpha x}.
\]
This convergence is uniform over \(z\) in a compact set. One can obtain
uniformity from convergence on a finite grid of interval endpoints and
monotone upper/lower interval approximations. The original entry is in a
compact interval, and the exit overshoot lies in \([0,y_+]\), where
\(y_+\) bounds the positive increments. Thus the averaging is valid
without any unbounded-entry interchange. The strong Markov property gives
the exact killed-potential identity
\[
 U_Z(D)=\mathbb E\sum_{n<\tau}\mathbf1_{\{Q_n\in D\}}
       +\mathbb E U_{Q_\tau}(D).
\]
Subtract the two uniformly valid asymptotics and divide by
\(\mathbb E\tau\) to obtain \eqref{sc:regeneration-asymptotic}.
Bounded local renewal masses give a uniform exponential bound for the
fixed-width intervals; a geometric-series domination then permits
summing toward \(-\infty\). Hence for \(R=e^Q\),
\[
 \mathbb P_\pi(R\le r)\sim
 \frac{D_\alpha}{\alpha m_\alpha\mathbb E\tau}r^\alpha.
\]
The bounded continuous increment density also verifies the density
version here directly. The tilted unrestricted potential satisfies
\(U_\alpha=\delta_0+F_\alpha*U_\alpha\). Its absolutely continuous
part therefore has the continuous density
\(u_\alpha(x)=\int f_\alpha(x-y)U_\alpha(dy)\).
The fixed-interval renewal theorem gives vague convergence of
\(U_\alpha(x+dt)\) to \(dt/m_\alpha\) as \(x\to-\infty\).
Because \(f_\alpha\) is continuous and compactly supported, this
convolution converges to \(1/m_\alpha\). The convergence is uniform
under bounded shifts. At sufficiently negative locations the compactly
supported entry and exit measures have no atoms; the killed-potential
subtraction therefore has a density there. Undoing the tilt and
subtracting its entry and exit coefficients proves
\(p_R(r)\sim D_\alpha r^{\alpha-1}/(m_\alpha\mathbb E\tau)\).
This result proves an excursion law with specified return. Applying it
to a nonlinear SGD path requires verifying its return and approximation
hypotheses, including the distribution of entry near zero.

\subsection{Proof of the smooth nonlinear SGD example}
For the general parameters in Theorem~\ref{sc:nonlinear-recurrence},
expansion of the sampled loss gives
\[
 \ell(x;h,e)=\tfrac12(h-\beta)x^2+\tfrac\beta2\log(1+x^2)
              -\sqrt h\,e x+\tfrac12e^2.
\]
Because \(h-\beta\ge\mu-s>0\), every realized loss is coercive.
Taking expectations gives
\(L(x)=\mu x^2/2+(\beta/2)\log(1+x^2)+C\).
Its derivative is \(x[\mu+\beta/(1+x^2)]\), so zero is the unique
minimizer. Its second derivative is
\[
 L''(x)=\mu+\beta\frac{1-x^2}{(1+x^2)^2}.
\]
The fraction attains its minimum \(-1/8\) at \(x^2=3\); thus
\(\min L''=\mu-\beta/8<0\).
Differentiation of the sampled loss gives the exact step
\[
 X^+=(1-\eta(h-\beta))X-\frac{\eta\beta X}{1+X^2}
            +\eta\sqrt h\,e.
\]
The remainder has absolute value at most \(\eta\beta/2\). The assumptions
on \(\eta\) imply
\(a_*:=\max\{|1-\eta(\mu-s)|,|1-\eta(\mu+s)|\}<1\).
For any \(p>0\), choose \(\delta>0\) so
\(q=(1+\delta)a_*^p<1\). The elementary bound
\[
 |u+v|^p\le(1+\delta)|u|^p+C_{p,\delta}|v|^p
\]
(or subadditivity when \(p\le1\)) gives
\(P|x|^p\le q|x|^p+C_p\), since the remainder is bounded and the
Gaussian noise has all moments. At \(p=1\) this drift, combined with the
Gaussian compact minorization in the proof of
Proposition~\ref{sc:outer-sign}, gives recurrence, uniqueness, and
convergence. This direct argument also covers a zero outer slope, which
was excluded from the logarithmic-sign proposition.
Iteration from any fixed finite start
bounds \(\sup_n\mathbb E|X_n|^p\). Convergence in law and the
Portmanteau theorem give \(\int|x|^p\pi(dx)<\infty\), without assuming
this moment when integrating the drift.

The conditional noise variances are \(\eta^2\sigma^2h\), bounded above
and away from zero. Invariance writes the density as a mixture of
translated Gaussian densities. Every Gaussian derivative of fixed
order is uniformly bounded over its translation and this compact
variance interval. Differentiating under the invariant mixture shows
that the density is bounded and smooth; positivity is immediate.
Reflection equivariance and uniqueness give evenness.
Finally the exact sampled Jacobian at zero is \(1-\eta h\).
At the numerical instance in the theorem, the signed gains are
\(-.91425,-1.11575\), whose product is greater than one.
Their harmonic mean is greater than one by direct rational arithmetic.
Strict convexity of \(q\mapsto\mathbb E|1-\eta h|^{-q}\), its negative
derivative at zero, and its divergence as \(q\to\infty\) (from the
contracting atom) give a unique positive root. Evaluating this finite
two-term function gives \eqref{sc:concrete-numbers}; the signs are exact
and do not depend on the displayed decimal rounding.

Increasing \(\beta\) while holding \((\mu,s,\eta)\) fixed leaves the
outer slopes unchanged and makes both central gains grow without bound.
Every fixed finite \(\beta\) still has a bounded remainder and finite
Gaussian moments. Consequently recurrence and all finite polynomial
moments can coexist with arbitrarily large central logarithmic and
harmonic gains, although a contracting central atom and hence a finite
inverse root eventually disappear.

\paragraph{A contrasting occupied inverse moment at a larger step.}
The same fixed loss at \(\eta=8\) has outer slopes \(.76,-.04\) and
central Jacobians \(-6.6,-7.4\), so all recurrence conclusions remain
valid and \(H_0=6.977142857\ldots>1\).

The Jacobian at an occupied state is instead
\[
 D_h(x)=1-8(h-.92)-7.36\frac{1-x^2}{(1+x^2)^2}.
\]
It is state dependent. In fact its simple zeros, together with the
strictly positive invariant density, imply
\(\mathbb E_{\pi,h}|D_h(X)|^{-1}=\infty\): the integrand is comparable
to \(1/|x-x_c|\) near any simple zero \(x_c\).
Thus the occupied harmonic gain can be zero while the central frozen
harmonic gain exceeds one. A finite Monte Carlo inverse moment need not
detect this divergence. This reinforces why all scalar directions and
locations must be identified in an empirical threshold claim.

\paragraph{When central depletion gives off-center modes.}
Suppose a signed invariant law has a continuous even density \(f\) on
\(\mathbb R\), \(f(0)=0\), \(f(x)\to0\) as \(|x|\to\infty\), and
positive total mass. Some \(x_0>0\) has \(f(x_0)>0\). The limit at zero
and infinity confines the supremum on \([0,\infty)\) to a compact
subinterval of \((0,\infty)\), where continuity attains it. Its negative
reflection is another maximizer. Thus there are at least two off-center
modes. This does not exclude additional modes or a plateau. Exactly two
requires, for example, a unique maximum on each half-line. Moreover,
at fixed nonzero additive Gaussian noise the density in the theorem is
positive at zero, so an ideal multiplicative depletion law cannot be
transferred literally down to zero without a small-noise argument.

\subsection{A counterexample to a universal harmonic-to-bimodality claim}
\label{sc:counterexample}
On \((0,1)\), take \(M=1/2\) with probability \(1/5\) and \(M=2\)
otherwise. If \(MR<1\), set \(R'=MR\); otherwise independently restart
from density
\[
 g(r)=\begin{cases}(24+21r^2)/41,&0<r<1/2,\\
                   (72-27r^2)/41,&1/2\le r<1.
       \end{cases}
\]
The multiplier has \(\gamma=(3/5)\log2>0\), \(H=5/4>1\), and inverse
root \(\alpha=2\). Nevertheless the exact invariant density is
\[
 f(r)=\frac{12}{11}(1-r^2/4),\qquad0<r<1,
\]
which decreases from the origin. Both displayed functions integrate to
one. The killed transfer operator is
\[
 Tf(r)=\frac25 f(2r)\mathbf1_{\{r<1/2\}}+\frac25 f(r/2).
\]
Direct substitution on the two half-intervals gives
\(f-Tf=(41/110)g\). Under \(f\), the reset probability is
\((4/5)\int_{1/2}^1f(r)dr=41/110\), so \(f=Tf+(41/110)g\)
proves invariance. Adding a reflection-symmetric sign yields a smooth
central maximum. The return law places substantial mass arbitrarily
near zero and does not satisfy the inverse entry-moment condition of
the more general regenerative theorem. The multiplier alone cannot
distinguish this counterexample from a regular-return model with the
same \(\alpha\).

\subsection{Exactly two modes in the confined diffusion normal form}
For clarity, one setting where mode count can be proved completely is
the \emph{diffusion} normal form
\[
 dY=(aY-cY^3)dt+\sqrt{\sigma^2+\nu^2Y^2}\,dW,
 \qquad c,\sigma,\nu>0.
\]
The quartic inward drift and nondegenerate diffusion give a nonexplosive
positive recurrent scalar process. Solving its zero-current equation
\((Dp)'=2(a y-c y^3)p\), where \(D=\sigma^2+\nu^2y^2\), gives
\[
 p(y)=C(\sigma^2+\nu^2y^2)^{a/\nu^2+c\sigma^2/\nu^4-1}
                e^{-cy^2/\nu^2}.
\]
It is normalizable, and direct differentiation yields
\[
 (\log p)'(y)=\frac{2y[(a-\nu^2)-cy^2]}{\sigma^2+\nu^2y^2}.
\]
There is one mode at zero for \(a\le\nu^2\), and exactly two modes at
\(\pm\sqrt{(a-\nu^2)/c}\) for \(a>\nu^2\).
The zero-additive-noise local logarithmic exponent is
\(a-\nu^2/2\). When $a>\nu^2/2$, its nontrivial zero-noise invariant density on the
punctured line has central exponent \(\alpha=2a/\nu^2-1\), so the mode transition occurs at
\(\alpha=1\). This is a diffusion theorem; fixed-step SGD requires its
own confinement and approximation argument. In particular, a raw cubic
explicit update with nondegenerate Gaussian innovations can escape and
does not inherit this invariant law merely because its local Taylor
expansion matches the diffusion drift.
